\section{Digital Music and the Rise of Home Production}\label{sec:digital-music-domain}

Recorded music is now a large digital ecosystem rather than a niche computing topic.
IFPI reported that global recorded music revenues reached US\$31.7 billion in 2025, with streaming representing 69.6\%
of total recorded music revenue and paid subscription streaming representing 52.4\% of the total
market~\cite{IFPI_2026}.
Spotify, reporting from the perspective of a single streaming platform, stated that it paid more than US\$11 billion to
the music industry in 2025~\cite{Spotify_2026_Payouts}.
BandLab's public product material describes a community of 100 million music creators using its platform
ecosystem~\cite{BandLab_2026}.
We cite these figures as industry-scale context, not as direct evidence of demand for a programming language.
They do show that digital music creation, distribution, and consumption are large enough for specialized creative
software to be a relevant engineering topic.

\subsection{From Commercial Studios to Personal Workstations}\label{subsec:home-production-shift}

Parallel to streaming growth, music production itself has decentralized.
Persson traces how home recording equipment and digital audio workstations created competition for commercial
studios and forced a re-evaluation of which production tasks happen at home, which remain in professional
facilities, and who performs them~\cite{Persson_2006}.
What began as project studios and hobby setups has become a sustained shift: affordable computers, audio
interfaces, and software replaced much of the dedicated hardware that once required a commercial facility.

Etinger surveyed 838 heavy-metal and hard-rock artists about DAW usage across composition, pre-production,
recording, mixing, and mastering~\cite{Etinger_2016}.
The study combined Task--Technology Fit with the Technology Acceptance Model and found that perceived usefulness
and ease of use predict sustained DAW adoption.
Together with Persson's industry analysis, this evidence supports a practical claim we rely on throughout the
thesis: for a growing population of creators, the primary composition environment is software on a personal
computer, not a score desk or a live-coding stage.

That shift does not eliminate professional studios, but it changes the default workflow.
Many composers now sketch, arrange, and iterate inside a DAW before any release-oriented mastering step.
Tools that interoperate with that workflow---especially through portable symbolic formats such as MIDI---remain
technically relevant even when the composer is not writing code.
